% >>> typage-template revtex-aps@2026.10.1 variant=prl
\documentclass[aps,prl,reprint,amsmath,amssymb]{revtex4-2}
% <<< typage-template
\usepackage{graphicx}
\usepackage{dcolumn}
\usepackage{bm}
\usepackage{booktabs}

\begin{document}

\preprint{APS/123-QED}

\title{Paper Title}

\author{First Author}
\email{first.author@example.com}
\affiliation{Department, University, City, Country}

\author{Second Author}
\affiliation{Institute, City, Country}

\date{\today}

\begin{abstract}
State the problem, the approach and the main result in a single
paragraph.
\end{abstract}

\maketitle

\section{Introduction}
Introduce the problem and list the contributions. Cite prior work by
number~\cite{lamport1994latex,knuth1984texbook}.

\section{Method}
Describe the method. Number the equations you refer to:
\begin{equation}
  \label{eq:objective}
  \min_{\theta} \; \frac{1}{n} \sum_{i=1}^{n} \ell\bigl(f_\theta(x_i), y_i\bigr) .
\end{equation}

\section{Results}
Table~\ref{tab:results} reports the results of
optimizing~\eqref{eq:objective}. Describe the data, the baselines and the
compute used, and report the variation across runs.

\begin{table}[!htb]
  \caption{Results on the example benchmark (mean $\pm$ standard deviation over five runs).}
  \label{tab:results}
  \centering
  \begin{tabular}{lcc}
    \toprule
    Method   & Accuracy (\%)  & Time (s) \\
    \midrule
    Baseline & 81.2 $\pm$ 0.4 & 12.4     \\
    Ours     & 84.7 $\pm$ 0.3 & 10.9     \\
    \bottomrule
  \end{tabular}
\end{table}

\section{Conclusion}
Conclude and outline future work.

\begin{acknowledgments}
Acknowledge funding and help here.
\end{acknowledgments}

\bibliography{refs}

\end{document}
